\chapter{Methodology and Requirements}

\section{Methodology and Toolchain Selection}
In physically constructing the HomeSOC application, a rapid, iterative agile development methodology was actively integrated. The overarching project timeline was aggressively sequenced into distinct, segmented, programmable objectives, meticulously prioritizing the core local data processing loop prior to engineering the eventual full-scale cloud deployment mechanism. This structural phase formally specifies the critical path toolchain selected for conceptualizing and finalizing the analytical system architecture.

\subsection{Back-End Programming Languages}
The foundational logic processing, ingestion API scripting, and physical agent polling sequences were entirely formulated utilising Python 3 (specifically interpreted under Python 3.8+ versions). This deliberate selection was primarily necessitated by Python's unparalleled ecosystem of mature, cross-platform external libraries combined with its native syntactical brevity, which enables rapid prototyping of computational architectures. 

Unlike heavy, compiled binaries derived from C or Go, which require specific architectural build environments, Python's ubiquitous interpreter nature ensures maximum logistical distribution across heterogeneous network spaces (running invariably on legacy Windows machines, UNIX-based Linux distributions, and containerised micro-services). The built-in `json` string encoder mapped telemetry flawlessly without requiring third-party serializers. Furthermore, the reliance on the third-party `psutil` (Process and System Utilities) library facilitated seamless, hardware-agnostic metric scraping of deep systemic constants (e.g., CPU frequency thresholds, active PID enumeration) isolated entirely from native operating systems shells.

\begin{figure}[h]
    \centering
    % Placeholder for Agent Output Screenshot
    \fbox{
        \begin{minipage}{0.6\textwidth}
            \vspace{2cm}
            \begin{center}
                \textit{[PLACEHOLDER: Insert Terminal Output Screenshot of the Python Host Agent actively scraping telemetry metrics (e.g. `python agent.py` command line output)]} \\
                \vspace{0.5cm}
                File: \texttt{agent\_console\_log.png}
            \end{center}
            \vspace{2cm}
        \end{minipage}
    }
    \caption{Expected native console output from the deployed telemetry agent operating successfully on a host node.}
    \label{fig:agent_console}
\end{figure}

\subsection{Microservices Framework Integration}
To accurately synthesize and bridge the localized agent traffic to the centralised backend, the Flask micro-framework routing logic was heavily utilized. This selection was diametrically opposed to heavier web administration paradigms like the Django MVC system, which explicitly imposes rigid monolithic configurations onto the project structure. 

Flask provided the absolute minimalistic flexibility essential for writing ultra-fast, highly targeted REST APIs capable of concurrently ingesting tens of thousands of localized TCP/IP telemetry strings devoid of excess computational bloat. 

\subsection{Data Persistence and Storage Mediums}
Data storage persistence natively relied on a robust, localized, serverless SQLite (Structured Query Language) installation. The conscious implementation decision to systematically abstain from deploying a sprawling relational database management system (RDBMS) like MySQL or a Document-Store NoSQL engine like MongoDB was principally dictated by the project's inviolable core tenet: extreme accessibility. The HomeSOC architecture had to absolutely remain logistically undemanding and universally executable across low-powered or legacy hardware environments quintessential to home lab networks.

However, as the project scaled from local testing to concurrent testing, the database inevitably confronted severe performance barriers and "I/O lock" anomalies when positioned within a simulated cloud filesystem (specifically, encountering heavy Network File System (NFS) input overhead queues). The root architecture was subsequently refined by adopting SQLite's Write-Ahead Logging (Pragmatic parameter `journal\_mode=WAL`) overriding the default legacy rollback journaling system to optimize multithreaded ingress.

\subsection{Front-End Aesthetics and UI Frameworks}
The aesthetic ambition behind the dashboard was the synthesis of a premium, corporate-grade graphical user interface (GUI) coupled with non-blocking, zero-latency visual polling. Consequently, standard, non-compiled web technologies dominated the UI specification. 

Vanilla ECMAScript JavaScript provided instantaneous Document Object Model (DOM) manipulations directly via the browser's asynchronous `fetch` integration method API. The UI eschews framework build steps (React, Node, Webpack). For styling rendering, complex CSS3 grid mechanics, specifically adopting the modern OKLCH hue parameter colour space calculation system, delivered robust dark-mode aesthetics that facilitate the projection of dynamic scalable interface elements natively.

Ultimately, achieving functional redundancy, the total source code repository underwent a critical final migration from local `localhost` network adapters (127.0.0.1 loopbacks) out into a full-scale production-grade cloud container infrastructure, orchestrated exclusively through the PythonAnywhere web-hosting utility provider.

\section{Software and Hardware Requirement Analysis}
Prior to programmatic execution, exhaustive conceptual and physical analyses establishing specific hardware dependencies and software interpreter prerequisites were systematically finalized to guarantee overarching viability metrics when scaling across diverse client networks.

\subsection{Endpoint Hardware Specifications}
A distinguishing novelty of the proposed local SOC application compared to competing enterprise products is its sheer nominal resource footprint. The requirements actively define computational scarcity:

\begin{enumerate}
    \item \textbf{Monitored Agent Endpoints:} The nodes can encompass any physical or hypervisor-virtual machine actively running a modern kernel operating system (distros ranging from Debian/Ubuntu Linux, Microsoft Windows 10/11, or macOS Big Sur upwards). The node mandates a restrictive minimum of merely 256MB allocated RAM and an operational standard Network Interface Card (NIC). No dedicated compute logic clusters (GPUs) are required, firmly placing the agent parameters within the capabilities of hardware as inexpensive and minimal as an ARM-based Raspberry Pi.
    \item \textbf{Analytical Central Server (Hybrid/Static):} When hosting the backend statically on physical hardware, the machine merely requires the presence of a baseline Consumer Dual-Core Processor (x86 or ARM64 instruction set), 1GB of continuous functional system memory, and a bare minimum of 5GB of localized storage capacity strictly dedicated to preserving historical tabular database logging records (`soc\_db.sqlite`). A static IP lease and continuous stable routing are explicitly crucial.
\end{enumerate}

\subsection{Software Stack Requirements}
Given the heavy dependence on OS-agnostic frameworks, the programmatic requisites are deliberately universalized:
\begin{enumerate}
    \item \textbf{Scripting Interpreter:} Python version 3.8, natively required in order to support executing precise f-string format literals, syntax iteration functions in tuples, and concurrent library modules deployed across the ingestion routing sequences.
    \item \textbf{Network Dependency Tree:} A virtualized shell environment `venv` physically retaining the package library list documented inside `requirements.txt`. Essential installation prerequisites included `Flask`, `Flask-Cors` (resolving complex cross-origin resource sharing blockages), `psutil`, `requests`, and potentially `mysql-connector-python` as a redundancy mechanism.
    \item \textbf{Analyst Client Display Access:} A highly standardized WebKit or Chromium-based modern visual browser.
    \item \textbf{Application Deployment Environment:} A functional PythonAnywhere WSGI configuration web application instance preconfigured over Ubuntu 22.04 LTS equipped natively with accessible secure bash consoles.
\end{enumerate}

\section{Methodology Table Synthesization}
The final toolchain assembly mapping represented systematically in Table 4-1 accurately reflects the strategic prioritization of security insight visibility combined with drastic systemic simplification over abstract operational complexity.

\begin{table}[h]
\centering
\small
\begin{tabularx}{\textwidth}{|X|X|X|}
\hline
\textbf{Architectural Category} & \textbf{Tool / Library Utilized} & \textbf{Strategic Selection Justification Matrix} \\ \hline
Backend API Framework & Python / Flask & Zero bloat, maximal REST integration flexibility, JSON ease. \\ \hline
System Telemetry Extractor & psutil Module & Hardware-agnostic deep systemic kernel scraping logic. \\ \hline
Data Persistence Archival & SQLite3 (WAL pragmas) & Total zero-configuration, standalone executable serverless storage format. \\ \hline
Production Cloud Infrastructure & PythonAnywhere & Persistent, universally accessible TLS terminal presence. \\ \hline
Frontend Interface Logic & Pure Vanilla JS / CSS3 & Infinite backward DOM compatibility array producing sub-millisecond reaction speeds. \\ \hline
\end{tabularx}
\caption{Comprehensive breakdown of the fundamental technology stack parameters dictating the HomeSOC framework construction.}
\label{tab:method_stack}
\end{table}

\section{Verification Specifications - Bruteforce Mathematical Analytical Plan}
Analytical project milestones heavily dictated that the completed ingestion API must faultlessly accept variable JSON payloads without dropping packets due to HTTP payload congestion overload. The system essentially must guarantee returning a standard `HTTP 200 OK` succession code or a strict `HTTP 400 Bad Request` verification error immediately if fundamental structural parsing mechanisms internally failed. 

To formally verify design adherence against these rigorous specifications, the testing phase integrated a dedicated, heavily multithreaded Python stress-testing script module specifically coded as `stress\_test.py`. This script was purposefully and aggressively utilized to brute-force the localized central API endpoint by computationally simulating an artificial multitude of heterogeneous nodes launching overlapping, consecutive mock attacks. Simulated parameters included generating massive synthesized CPU frequency throttling warnings, creating dummy SSH authentication login sequence denials, and embedding fake network footprint scanning payloads containing literal string variables commonly mapping to hostile binaries (`xmrig` cryptocurrency executable patterns and `nmap` reconnaissance routines). 

The defined verification criteria demanded compliance metrics indicating that:
\begin{itemize}
    \item The relational schema correctly mapped and appended 100\% of the simulated telemetry traffic logs seamlessly devoid of connection reset flags.
    \item The fundamental `detection.py` logical arithmetic engine correctly parsed JSON dictionaries and flawlessly aligned simulated mathematical telemetry values successfully against static pre-established numerical algorithmic thresholds encoded into the rules table.
    \item The primary analyst reporting dashboard effectively polled the resultant changes sequentially over network parameters and visibly assigned physical styling classifications of `Low`, `Medium`, and `High` severity tags to the specific anomalies accurately without missing data arrays. 
    \item The GUI interactive notification alert system logically and systematically strictly observed distinct threshold gating logic patterns (e.g. perpetually silencing modal pop-up alarms for minimal `Low` warnings, effectively constraining visual noise, whilst immediately and aggressively asserting graphical dominance and stacking arrays to demand interactive human operator acknowledgment strictly for prioritized `High` warnings).
\end{itemize}
This precise verification specification structure functionally informed the subsequent complex programmatic implementation mapping documented precisely in Chapter 5.

\section{Expanded Theoretical Framework: Software Development Lifecycles}
A software development process prescribes a process for developing software. It typically divides an overall effort into smaller steps or sub-processes that are intended to ensure high-quality results. The process may describe specific deliverables – artifacts to be created and completed. 
Although not strictly limited to it, software development process often refers to the high-level process that governs the development of a software system from its beginning to its end of life – known as a methodology, model or framework. The system development life cycle (SDLC) describes the typical phases that a development effort goes through from the beginning to the end of life for a system – including a software system. A methodology prescribes how engineers go about their work in order to move the system through its life cycle. A methodology is a classification of processes or a blueprint for a process that is devised for the SDLC. For example, many processes can be classified as a spiral model.
Software process and software quality are closely interrelated; some unexpected facets and effects have been observed in practice.


== Methodology ==
The SDLC drives the definition of a methodology in that a methodology must address the phases of the SDLC. Generally, a methodology is designed to result in a high-quality system that meets or exceeds expectations (requirements) and is delivered on time and within budget even though computer systems can be complex and integrate disparate components. Various methodologies have been devised, including waterfall, spiral, agile, rapid prototyping, incremental, and synchronize and stabilize.
A major difference between methodologies is the degree to which the phases are sequential vs. iterative. Agile methodologies, such as XP and scrum, focus on lightweight processes that allow for rapid changes. Iterative methodologies, such as Rational Unified Process and dynamic systems development method, focus on stabilizing project scope and iteratively expanding or improving products. Sequential or big-design-up-front (BDUF) models, such as waterfall, focus on complete and correct planning to guide larger projects and limit risks to successful and predictable results. Anamorphic development is guided by project scope and adaptive iterations. In scrum, for example, one could say a single user story goes through all the phases of the SDLC within a two-week sprint. By contrast the waterfall methodology, where every business requirement is translated into feature/functional descriptions which are then all implemented typically over a period of months or longer.
A project can include both a project life cycle (PLC) and an SDLC, which describe different activities. According to Taylor (2004), "the project life cycle encompasses all the activities of the project, while the systems development life cycle focuses on realizing the product requirements".


=== History ===
The term SDLC is often used as an abbreviated version of SDLC methodology. Further, some use SDLC and traditional SDLC to mean the waterfall methodology.
According to Elliott (2004), SDLC "originated in the 1960s, to develop large scale functional business systems in an age of large scale business conglomerates. Information systems activities revolved around heavy data processing and number crunching routines". The structured systems analysis and design method (SSADM) was produced for the UK government Office of Government Commerce in the 1980s. Ever since, according to Elliott (2004), "the traditional life cycle approaches to systems development have been increasingly replaced with alternative approaches and frameworks, which attempted to overcome some of the inherent deficiencies of the traditional SDLC". The main idea of the SDLC has been "to pursue the development of information systems in a very deliberate, structured and methodical way, requiring each stage of the life cycle––from the inception of the idea to delivery of the final system––to be carried out rigidly and sequentially" within the context of the framework being applied.
Other methodologies were devised later:

1970s
Structured programming since 1969
Cap Gemini SDM, originally from PANDATA, the first English translation was published in 1974. SDM stands for System Development Methodology
1980s
Structured systems analysis and design method (SSADM) from 1980 onwards
Information Requirement Analysis/Soft systems methodology
1990s
Object-oriented programming (OOP) developed in the early 1960s and became a dominant programming approach during the mid-1990s
Rapid application development (RAD), since 1991
Dynamic systems development method (DSDM), since 1994
Scrum, since 1995
Team software process, since 1998
Rational Unified Process (RUP), maintained by IBM since 1998
Extreme programming, since 1999
2000s
Agile Unified Process (AUP) maintained since 2005 by Scott Ambler
Disciplined agile delivery (DAD) Supersedes AUP
2010s
Scaled Agile Framework (SAFe)
Large-Scale Scrum (LeSS)
DevOps
Since DSDM in 1994, all of the methodologies on the above list except RUP have been agile methodologies - yet many organizations, especially governments, still use pre-agile processes (often waterfall or similar). 


=== Examples ===
The following are notable methodologies somewhat ordered by popularity.

Agile
Agile software development refers to a group of frameworks based on iterative development, where requirements and solutions evolve via collaboration between self-organizing cross-functional teams. The term was coined in the year 2001 when the Agile Manifesto was formulated.

Waterfall
The waterfall model is a sequential development approach, in which development flows one-way (like a waterfall) through the SDLC phases.

Spiral
In 1988, Barry Boehm published a software system development spiral model, which combines key aspects of the waterfall model and  rapid prototyping, in an effort to combine advantages of  top-down and bottom-up concepts. It emphases a key area many felt had been neglected by other methodologies: deliberate iterative risk analysis, particularly suited to large-scale complex systems.

Incremental
Various methods combine linear and iterative methodologies, with the primary objective of reducing inherent project risk by breaking a project into smaller segments and providing more ease-of-change during the development process.

Prototyping
Software prototyping is about creating prototypes, i.e. incomplete versions of the software program being developed.

Rapid
Rapid application development (RAD) is a methodology which favors iterative development and the rapid construction of prototypes instead of large amounts of up-front planning. The "planning" of software developed using RAD is interleaved with writing the software itself. The lack of extensive pre-planning generally allows software to be written much faster and makes it easier to change requirements.

Shape Up
Shape Up is a software development approach introduced by Basecamp in 2018. It is a set of principles and techniques that Basecamp developed internally to overcome the problem of projects dragging on with no clear end. Its primary target audience is remote teams. Shape Up has no estimation and velocity tracking, backlogs, or sprints, unlike waterfall, agile, or scrum. Instead, those concepts are replaced with appetite, betting, and cycles. As of 2022, besides Basecamp, notable organizations that have adopted Shape Up include UserVoice and Block.

Chaos
Chaos model has one main rule: always resolve the most important issue first.

Incremental funding
Incremental funding methodology - an iterative approach.

Lightweight
Lightweight methodology - a general term for methods that only have a few rules and practices.

Agile software development is an umbrella term for approaches to developing software that reflect the values and principles agreed upon by The Agile Alliance, a group of 17 software practitioners, in 2001. As documented in their Manifesto for Agile Software Development, the practitioners value: 

Individuals and interactions over processes and tools
Working software over comprehensive documentation
Customer collaboration over contract negotiation
Responding to change over following a plan
The practitioners cite inspiration from new practices at the time including extreme programming, scrum, dynamic systems development method, adaptive software development, and being sympathetic to the need for an alternative to documentation-driven, heavyweight software development processes.
Many software development practices emerged from the agile mindset. These agile-based practices, sometimes called Agile (with a capital A), include requirements, discovery, and solutions improvement through the collaborative effort of self-organizing and cross-functional teams with their customer(s)/end user(s).
While there is much anecdotal evidence that the agile mindset and agile-based practices improve the software development process, the empirical evidence is limited and less than conclusive.


== History ==
Iterative and incremental software development methods can be traced back as early as 1957, with evolutionary project management and adaptive software development emerging in the early 1970s.
During the 1990s, a number of lightweight software development methods evolved in reaction to the prevailing heavyweight methods (often referred to collectively as waterfall) that critics described as overly regulated, planned, and micromanaged. These lightweight methods included rapid application development (RAD), from 1991; the unified process (UP) and dynamic systems development method (DSDM), both from 1994; Scrum, from 1995; Crystal Clear and extreme programming (XP), both from 1996; and feature-driven development (FDD), from 1997. Although these all originated before the publication of the Agile Manifesto, they are now collectively referred to as agile software development methods.
Already since 1991 similar changes had been underway in manufacturing and management thinking derived from lean management.
In 2001, seventeen software developers met at a resort in Snowbird, Utah to discuss lightweight development methods. They were Kent Beck (Extreme Programming), Ward Cunningham (Extreme Programming), Dave Thomas (Pragmatic Programming, Ruby), Jeff Sutherland (Scrum), Ken Schwaber (Scrum), Jim Highsmith (Adaptive Software Development), Alistair Cockburn (Crystal), Robert C. Martin (SOLID), Mike Beedle (Scrum), Arie van Bennekum, Martin Fowler (OOAD and UML), James Grenning, Andrew Hunt (Pragmatic Programming, Ruby), Ron Jeffries (Extreme Programming), Jon Kern, Brian Marick (Ruby, Test-driven development), and Steve Mellor (OOA). The group, The Agile Alliance, published the Manifesto for Agile Software Development.
In 2005, a group headed by Cockburn and Highsmith wrote an addendum of project management principles, the PM Declaration of Interdependence, to guide software project management according to agile software development methods.
In 2009, a group working with Martin wrote an extension of software development principles, the Software Craftsmanship Manifesto, to guide agile software development according to professional conduct and mastery.
In 2011, the Agile Alliance created the Guide to Agile Practices (renamed the Agile Glossary in 2016), an evolving open-source compendium of the working definitions of agile practices, terms, and elements, along with interpretations and experience guidelines from the worldwide community of agile practitioners.


== Values and principles ==


=== Values ===
The manifesto for agile software development reads:
We are uncovering better ways of developing software by doing it and helping others do it. Through this work we have come to value:

  Individuals and interactions over processes and tools 
  Working software over comprehensive documentation 
  Customer collaboration over contract negotiation 
  Responding to change over following a plan 
That is, while there is value in the items on the right, we value the items on the left more.
Scott Ambler explained:

Tools and processes are important, but it is more important to have competent people working together effectively.
Good documentation is useful in helping people to understand how the software is built and how to use it, but the main point of development is to create software, not documentation.
A contract is important but is not a substitute for working closely with customers to discover what they need.
A project plan is important, but it must not be too rigid to accommodate changes in technology or the environment, stakeholders' priorities, and people's understanding of the problem and its solution.
Introducing the manifesto on behalf of the Agile Alliance, Jim Highsmith said,

The Agile movement is not anti-methodology, in fact many of us want to restore credibility to the word methodology. We want to restore a balance. We embrace modeling, but not in order to file some diagram in a dusty corporate repository. We embrace documentation, but not hundreds of pages of never-maintained and rarely-used tomes. We plan, but recognize the limits of planning in a turbulent environment. Those who would brand proponents of XP or SCRUM or any of the other Agile Methodologies as "hackers" are ignorant of both the methodologies and the original definition of the term hacker.


=== Principles ===
The values are based on these principles:

Customer satisfaction by early and continuous delivery of valuable software.
Welcome changing requirements, even in late development.
Deliver working software frequently (weeks rather than months).
Close, daily cooperation between business people and developers.
Projects are built around motivated individuals, who should be trusted.
Face-to-face conversation is the best form of communication (co-location).
Working software is the primary measure of progress.
Sustainable development, able to maintain a constant pace.
Continuous attention to technical excellence and good design.
Simplicity—the art of maximizing the amount of work not done—is essential.
Best architectures, requirements, and designs emerge from self-organizing teams.
Regularly, the team reflects on how to become more effective, and adjusts accordingly.


== Overview ==


=== Iterative, incremental, and evolutionary ===
Most agile development methods break product development work into small increments that minimize the amount of up-front planning and design. Iterations, or sprints, are short time frames (timeboxes) that typically last from one to four weeks. Each iteration involves a cross-functional team working in all functions: planning, analysis, design, coding, unit testing, and acceptance testing. At the end of the iteration a working product is demonstrated to stakeholders. This minimizes overall risk and allows the product to adapt to changes quickly. An iteration might not add enough functionality to warrant a market release, but the goal is to have an available release (with minimal bugs) at the end of each iteration. Through incremental development, products have room to "fail often and early" throughout each iterative phase instead of drastically on a final release date. Multiple iterations might be required to release a product or new features. Working software is the primary measure of progress.


=== Efficient and face-to-face communication ===
The 6th principle of the agile manifesto for software development states, "The most efficient and effective method of conveying information to and within a development team is face-to-face conversation". The manifesto, written in 2001 when video conferencing was not widely used, states this in relation to the communication of information, not necessarily that a team should be co-located.
The principle of co-location is that co-workers on the same team should be situated together to better establish the identity as a team and to improve communication. This enables face-to-face interaction, ideally in front of a whiteboard, that reduces the cycle time typically taken when questions and answers are mediated through phone, persistent chat, wiki, or email. With the widespread adoption of remote working during the COVID-19 pandemic and changes to tooling, more studies have been conducted around co-location and distributed working which show that co-location is increasingly less relevant.
No matter which development method is followed, every team should include a customer representative (known as product owner in Scrum). This representative is agreed by stakeholders to act on their behalf and makes a personal commitment to being available for developers to answer questions throughout the iteration. At the end of each iteration, the project stakeholders together with the customer representative review progress and re-evaluate priorities with a view to optimizing the return on investment (ROI) and ensuring alignment with customer needs and company goals. The importance of stakeholder satisfaction, detailed by frequent interaction and review at the end of each phase, is why the approach is often denoted as a customer-centered methodology.


==== Information radiator ====
In agile software development, an information radiator is a (normally large) physical display, board with sticky notes or similar, located prominently near the development team, where passers-by can see it. It presents an up-to-date summary of the product development status. A build light indicator may also be used to inform a team about the current status of their product development.


=== Very short feedback loop and adaptation cycle ===
A common characteristic in agile software development is the daily stand-up (known as daily scrum in the Scrum framework). In a brief session (e.g., 15 minutes), team members review collectively how they are progressing toward their goal and agree whether they need to adapt their approach. To keep to the agreed time limit, teams often use simple coded questions (such as what they completed the previous day, what they aim to complete that day, and whether there are any impediments or risks to progress), and delay detailed discussions and problem resolution until after the stand-up.
